\chapter{Protocolo de Reproducibilidad}
\label{ap:reproducibilidad}

Este apéndice documenta el protocolo necesario para reproducir los resultados reportados en la tesis. La reproducción completa requiere el código experimental, los conjuntos de datos públicos indicados en el Capítulo~\ref{ch:metodologia}, las versiones de bibliotecas registradas en los archivos de configuración y los artefactos derivados conservados junto con el documento.

\section{Entorno computacional}

El sistema experimental fue desarrollado en Python y utiliza bibliotecas científicas estándar: MNE-Python para manejo de registros EEG, NumPy y SciPy para álgebra numérica, scikit-learn para métodos auxiliares, Optuna para optimización de hiperparámetros, pandas para manejo de tablas y matplotlib para figuras. Las versiones exactas deben fijarse antes de ejecutar la reproducción porque cambios menores en solvers, interpoladores o generadores de números aleatorios pueden modificar resultados de frontera.

\section{Datos y preprocesamiento}

Los conjuntos de datos se obtienen desde fuentes públicas. Cada registro se transforma a una interfaz común: matriz de canales por tiempo, nombres de canales, frecuencia de muestreo, coordenadas de electrodos y metadatos de sujeto o ensayo. El preprocesamiento se mantiene idéntico entre métodos: filtrado, referencia, segmentación y selección de canales se aplican antes de simular pérdida de canales. Cuando un canal no dispone de coordenadas confiables, se excluye de afirmaciones fuertes.

\section{Máscaras y semillas}

La pérdida de canales se simula mediante máscaras aleatorias y contiguas. Cada máscara queda definida por conjunto de datos, severidad, modo de pérdida y semilla. La misma máscara se aplica a todos los métodos de un caso pareado, lo que permite atribuir diferencias de desempeño al método y no a variaciones del caso experimental. Las semillas deben fijarse globalmente antes de generar máscaras, optimizar hiperparámetros o seleccionar casos representativos.

\section{Fases reproducibles}

La reproducción sigue cuatro fases. Primero se ejecuta la selección preliminar de métodos y grafos, que explora familias amplias sin convertir sus rankings en evidencia final. Luego se realiza la reducción de candidatos, documentando costo, circularidad y reproducibilidad. En tercer lugar se ejecuta la optimización intermedia y se congelan variantes. Finalmente se corre la comparación pareada TRSS--MNE y se generan tablas/figuras del Capítulo~\ref{ch:resultados}.

\section{Reglas de exclusión}

Las reglas de exclusión deben aplicarse antes de sintetizar resultados. Se descartan ejecuciones con valores no numéricos, tiempos excesivos, identificadores duplicados inconsistentes, reconstrucciones triviales idénticas a la entrada o variantes que no formen parte del conjunto final de métodos válidos para análisis. Solo después de aplicar esos filtros se generan rankings exploratorios, tablas descriptivas y figuras del documento.

\section{Validación del documento}

La versión final del documento debe compilarse desde el controlador principal y, cuando corresponda, desde controladores parciales usados para revisión. La validación mínima incluye: ausencia de errores LaTeX, ausencia de citas o referencias indefinidas, ausencia de overfull hbox, comprobación de figuras existentes, y revisión visual de páginas con tablas/figuras modificadas. Las tablas del Capítulo~\ref{ch:resultados}, Capítulo~\ref{ch:conclusiones} y anexos se generan o verifican contra los CSV derivados antes de cerrar el documento.
